\subsection{AC-G86: an unconditional explicit bound for the actual auxiliary prime}
\label{sec:cubic-effective-actual-prime}

\paragraph{Scope.}
This section uses the actual branch and scalar element constructed in AC-G78 and AC-G83, together with the fixed-precision conclusion of AC-G84 and the finite order quotient of AC-G85. The conclusions of AC-G84 and AC-G85 are used with their checked normalization variant: trace congruences are required in the normalized algebra \(S\). The result bounds the prime in this actual construction. It does not establish the polynomial budget required by the inherited counting argument or containment in its ideal \(J\). No Riemann hypothesis for a nonabelian splitting field is inferred from A.

\paragraph{Notation and published input.}
Put \(p=269\). Let \(\mathcal O\) be the ring of integers of a finite extension \(E_0/\mathbb Q_p\), put \(c=[E_0:\mathbb Q_p]\), and let \(B=\mathcal O[[W]]\). Let \(S\) be the normalization of the selected irreducible Hida branch, finite free of rank \(D\) over \(B\). Thus AC-G85 gives a continuous finite quotient of \(G_{\mathbb Q}\), unramified outside the source ramification set, of order at most
\[
 p^{64Dc},
\]
such that equality in this quotient implies equality of the actual branch trace modulo
\(I S\), where \(I=(p^4,W^4)\subset B\).
Let \(\Sigma\) be a finite set of rational primes containing \(2,3,p\), every prime of the source tame level, every ramified prime of the quadratic hosts and of the real quadratic field \(E\) used in AC-G83, and any further finite set of primes which must be excluded from the choice of \(q\). Define
\[
 P_\Sigma=\prod_{\ell\in\Sigma}\ell,\qquad
 m_\Sigma=8\prod_{\substack{\ell\in\Sigma\\\ell\ {\rm odd}}}\ell,
 \qquad
 N_*=48m_\Sigma(p-1)^2p^{64Dc+4}.
\]
The set \(\Sigma\), the ring \(S\), and \(E\) are fixed before \(q\) is chosen.

The imported prime theorem is Kadiri--Wong, \emph{Primes in the Chebotarev density theorem for all number fields}, J.\ Number Theory 241 (2022), 700--737, Theorem 1, printed p.\ 701. For a nontrivial finite Galois extension \(\mathcal L/\mathbb Q\), each conjugacy class has an unramified rational prime \(q\leq |\operatorname{Disc}(\mathcal L)|^{310}\). This is an unconditional theorem for all discriminants. Its published numerical verification in Appendix A is retained as part of that imported result; no numerical rerun is claimed here.

\paragraph{Finding AC-G86-disc: a discriminant bound from degree and ramification.}
If \(\mathcal L/\mathbb Q\) is finite Galois of degree \(n\) and unramified outside \(\Sigma\), then
\[
 |\operatorname{Disc}(\mathcal L)|\leq (nP_\Sigma)^n. \tag{G86.1}
\]
Here is a proof including the wild-prime estimate. Fix \(\ell\in\Sigma\), a completion \(L_\ell\), its ramification index \(e\), and the maximal unramified subfield \(U/\mathbb Q_\ell\) of \(L_\ell/\mathbb Q_\ell\). The extension \(L_\ell/U\) is totally ramified of degree \(e\). Choose a uniformizer \(\varpi\) of \(L_\ell\). Its minimal polynomial
\[
 f(X)=X^e+a_{e-1}X^{e-1}+\cdots+a_0
\]
over \(U\) is Eisenstein and \(\mathcal O_{L_\ell}=\mathcal O_U[\varpi]\). These facts also follow directly from valuations: the powers \(1,\varpi,\ldots,\varpi^{e-1}\) have different valuations modulo \(e\), so they are a basis and integrality of a linear combination forces each coefficient to be integral. The usual monogenic trace-dual calculation gives the different as \(f'(\varpi)\). To recall that calculation, Lagrange interpolation gives
\(\operatorname{Tr}(\varpi^j/f'(\varpi))=0\) for \(0\leq j<e-1\) and \(1\) for \(j=e-1\); reducing higher powers with the monic polynomial makes the matrix \(\bigl(\operatorname{Tr}(\varpi^{i+j}/f'(\varpi))\bigr)_{0\leq i,j<e}\) integral, zero when \(i+j<e-1\), and equal to one when \(i+j=e-1\). Its determinant is therefore a unit. This proves that the trace-dual lattice is exactly \(f'(\varpi)^{-1}\mathcal O_{L_\ell}\).

The nonzero terms \(i a_i\varpi^{i-1}\), with \(a_e=1\), of \(f'(\varpi)\) have valuations congruent to \(i-1\) modulo \(e\). They are pairwise distinct modulo \(e\); consequently their minimum cannot cancel. In particular the different exponent \(d_\ell\), normalized on \(L_\ell\), satisfies
\[
 d_\ell\leq e-1+e\,v_\ell(e).
\]
The unramified extension \(U/\mathbb Q_\ell\) contributes no different. Since the global extension is Galois, all primes over \(\ell\) have these same local invariants and \(e\mid n\). Taking the norm of the different gives
\[
 v_\ell(\operatorname{Disc}(\mathcal L))
   ={n\over e}d_\ell
   \leq n\bigl(1+v_\ell(n)\bigr).
\]
Multiplication over \(\Sigma\) proves (G86.1), because
\(\prod_{\ell\in\Sigma}\ell^{v_\ell(n)}\leq n\).
The elementary local inputs are located, with the same valuation normalization, in Milne, \emph{Algebraic Number Theory}, v3.08, Proposition 7.55 and Remark 7.56, printed pp.\ 129--130; the trace-dual calculation above supplies the additional different formula used here.

\paragraph{Finding AC-G86: the actual prime can be chosen below an explicit bound.}
Under the preceding hypotheses the prime in AC-G78/AC-G83 can be chosen so that all the following hold simultaneously:
\[
 \begin{gathered}
 q\notin\Sigma,\qquad q\text{ splits completely in }F,\qquad
 \chi(q)=1,\qquad v_p(q+1)=1,\\
 q\text{ is inert in }E,\qquad
 u_E\in(\mathbb F_{q^2}^{\,\times})^p,\\
 T(\operatorname{Frob}^{\rm geom}_q)-(q+q^{-1})
       \text{ has an actual zero}\\
       \text{on the selected branch},\\
 q\leq (N_*P_\Sigma)^{310N_*}. \tag{G86.2}
 \end{gathered}
\]
Here \(u_E\) is the square of a fundamental unit as in AC-G83. The last zero is at a finite coefficient-field point with \(|W-p|<p^{-2}\); it is the actual Hecke specialization, with all the types retained in AC-G79 and AC-G80.

\emph{Proof.}
Let \(\mathcal L_0\) be the finite Galois field cut out by AC-G85's quotient. Form the compositum \(\mathcal L\) of \(\mathcal L_0\) with the following fields: the quartic host \(F\); the \(S_3\)-splitting field of \(\chi_{269,3}\), of degree \(6\); the field
\[
 M=E(\mu_p,u_E^{1/p})
\]
of AC-G83; \(\mathbb Q(\mu_{p^4})\); and \(\mathbb Q(\mu_{m_\Sigma})\).
The degree of \(M/\mathbb Q\) is at most \(2p(p-1)\), since its Kummer part has degree at most \(p\) over \(E(\mu_p)\). All these fields are Galois over \(\mathbb Q\). The elementary product bound for degrees yields
\[
 [\mathcal L:\mathbb Q]
 \leq p^{64Dc}\cdot4\cdot6\cdot2p(p-1)
          \cdot(p-1)p^3\cdot m_\Sigma
 =N_*.
\]
Every field in the compositum is unramified outside \(\Sigma\). For \(M\) this follows because \(u_E\) is a global unit: away from \(p\), the algebra obtained by adjoining its \(p\)-th root is finite etale, as the root and the derivative \(pX^{p-1}\) are units. The only additional ramification comes from \(E\) and \(\mu_p\), already included in \(\Sigma\). The field \(\mathbb Q(\mu_{m_\Sigma})\) is ramified at every prime in \(\Sigma\): its conductor has a factor \(8\) at \(2\) and a factor \(\ell\) at every odd \(\ell\in\Sigma\). Hence every unramified rational prime for \(\mathcal L\) automatically lies outside \(\Sigma\).

Let \(g=h't\in G_F\) be the actual conjugated scalar element of AC-G83. By construction
\[
 \varepsilon(g)=-(1+p),\quad
 \chi(g)=1,\quad
 \eta_E(g)=-1,\quad
 g^2|_M=1,
\]
and its actual branch trace is
\[
 T(g)=-\bigl((1+W)+(1+W)^{-1}\bigr).
\]
There is no independent compositum-compatibility assumption: \(g\) is already an element of the absolute Galois group. Apply the imported Chebotarev theorem to the conjugacy class of \(g^{-1}|_{\mathcal L}\), so that the resulting \emph{geometric} Frobenius has class \(g|_{\mathcal L}\). This convention matters, since
\(\varepsilon(\operatorname{Frob}^{\rm geom}_q)=q^{-1}\).

Equality in the finite order quotient gives
\(T(\operatorname{Frob}^{\rm geom}_q)-T(g)\in IS\).
Equality on \(\mu_{p^4}\) gives
\(q^{-1}\equiv-(1+p)\pmod{p^4}\). Inversion preserves this congruence of units, so the difference of the two cyclotomic sums also lies in \(p^4\mathcal O\). Consequently
\[
 \begin{split}
 F_q(W)&=T(\operatorname{Frob}^{\rm geom}_q)-(q+q^{-1}),\\
 F_g(W)&=-\bigl((1+W)+(1+W)^{-1}\bigr)
                  +(1+p)+(1+p)^{-1}
 \end{split}
 \qquad\text{satisfy}\qquad F_q-F_g\in IS .
\]
AC-G84 now supplies the actual zero. In particular, this argument uses the fixed quotient needed for that proof, not an unspecified neighborhood in a normalization.

The remaining finite field conditions give precisely the prime conditions in AC-G83: splitting in \(F\), trivial cubic Frobenius, inertness in \(E\), and order two on \(M\). Thus the Frobenius of the inert degree-two prime of \(E\) is trivial in \(M/E\); reduction gives \(u_E\) as a \(p\)-th power in \(\mathbb F_{q^2}\). The congruence modulo \(p^2\) gives \(v_p(q+1)=1\), since \(q^{-1}+1\equiv-p\pmod{p^2}\).
Finally, (G86.1), monotonicity of \(x\mapsto(xP_\Sigma)^x\) for \(x\geq1\), and the prime theorem give
\[
 q\leq|\operatorname{Disc}(\mathcal L)|^{310}
       \leq(N_*P_\Sigma)^{310N_*},
\]
as asserted. The actual representation and its ray-class cover then follow from the already proved constructions; no new automorphic existence assertion is hidden in the numerical step. \(\square\)

\paragraph{What the estimate settles and what it leaves.}
It supplies an unconditional, explicit upper bound for a prime satisfying the \emph{simultaneous actual} Hida trace and Kummer-cover conditions, in terms of the specified finite branch rank, coefficient degree, and ramification set. The bound is extremely large. It neither proves the inherited polynomial counting inequality nor places an entire curve in \(J\). In particular, different deeper approximation requests can still require different primes and different hosts. The Frame.Ok/EC/Big cohomological inequalities, the two-root replacements, and nonvanishing of the actual rank-two singular regulator remain separate gates.

\paragraph{Geometric dictionary.}
The finite continuous quotient in this proof corresponds to a finite etale Galois cover of \(\operatorname{Spec}\mathbb Z[1/P_\Sigma]\); adjoining the other fields takes the common refinement of these covers. Chebotarev chooses a closed point with the specified Frobenius conjugacy class on that refinement. This uses the already pinned SGA~1, Expose~V finite-cover/Galois-category dictionary. It does not identify the different characteristic-\(p\) automorphic support condition with a Frobenius conjugacy condition.
